\section{Spatial Discretization}
\label{sec:spatial_discretization}

\subsection{Discontinuous Galerkin Spectral Element Method}
\label{sec:dgsem}

The spatial \glspl{dgsem} discretizations developed in this section are based on the work of \cite{kronbichler_discontinuous_2021} and \cite{winters_construction_2021} and the book by \cite{hesthaven_nodal_2008}.

\subsubsection{Nodal Element Toolbox}
\label{sec:nodal_element_toolbox}

In the following, we construct matrix operators for polynomial interpolation, integration, differentiation, and integration by parts, as required by the \gls{dg} discretizations in \cref{sec:dg_discretization}.
To this end, we construct nodal \gls{dg} elements on the computational reference element $[-1,1]$ on Legendre-Gau\ss-Lobatto (LGL) points $\{\xi_k\}_{k=0}^{N^\mathrm{d}} \subset [-1,1]$ using Lagrange polynomials of degree $N^\mathrm{d}$ with basis $\mathcal{L}^{N^\mathrm{d}} \coloneq \{\ell_j\}_{j=0}^{N^\mathrm{d}}$,
\begin{align}
\label{eq:lagrange_basis_function}
\ell_j(\xi) = \prod_{\substack{i=0 \\ i \neq j}}^{N^\mathrm{d}} \frac{\xi-\xi_i}{\xi_j-\xi_i} \quad \text{such that} \quad \ell_j(\xi_i)= \delta_{ij} \coloneq 
\begin{cases}
1, & i=j,\\
0, & \text{else.}
\end{cases}
\end{align}
Consequently, the interpolatory polynomial $c_h\in \mathbb{P}^{N^\mathrm{d}}([-1,1])$ of a sought solution $c\in\mathcal{C}^\infty(E)$ takes the form
\begin{align}
c_h(\xi) &= \sum_{j=0}^{N^\mathrm{d}} \ell_j(\xi) c^{\mathrm{b}}(\xi_j) \approx c^{\mathrm{b}}(\xi).
\end{align}

Mass matrices are used to discretize integral operators and are defined as
\begin{align}
\label{eq:operator_mass_matrix}
\mathcal{M}_{i, j} &\coloneq \int_{-1}^{1} \ell_i(\xi) \ell_j(\xi) w(\xi) \,\mathrm{d}\xi \in\mathbb{R},
\end{align}
where $w\colon [-1,1]\to\mathbb{R}$ is a weighting function that, for example, accounts for geometric metric terms or spatially dependent parameters.
A particularly important mass matrix is obtained by choosing the Jacobi weight function $w^{(\alpha,\beta)} = (1-\xi)^\alpha (1+\xi)^\beta$ with $\alpha,\beta\in\mathbb{N}^0$, yielding
\begin{align}
\label{eq:operator_exact_jacobi_mass_matrix}
\mathcal{M}^{(\alpha, \beta)}_{i, j} &\coloneq \int_{-1}^{1} \ell_i(\xi) \ell_j(\xi) w^{(\alpha, \beta)}(\xi) \,\mathrm{d}\xi = \left(\mathcal{V}^{(\alpha, \beta)} \left( \mathcal{V}^{(\alpha, \beta)} \right)^T \right)^{-1}_{i,j},
\end{align}
where $\mathcal{V}^{(\alpha, \beta)}_{i, j} \coloneq \mathcal{P}^{(\alpha, \beta)}_j(\xi_i)$ defines the (invertible) Vandermonde matrix associated with the corresponding Jacobi polynomial basis $\left\{\mathcal{P}^{(\alpha, \beta)}_j\right\}_{j=0}^{N^\mathrm{d}}$. 
Alternatively, the weighted mass matrix can be approximated using numerical quadrature with $N^\mathrm{Q}\in\mathbb{N}$ quadrature points $\{\tilde{\xi}_k\}_{k=1}^{N^\mathrm{Q}}$,
\begin{align}
\label{eq:operator_gauss_mass_matrix}
\tilde{\mathcal{M}}_{i,j} &\coloneq \sum_{k=1}^{N^\mathrm{Q}} \ell_i(\tilde{\xi}_k) \ell_j(\tilde{\xi}_k) \, w(\tilde{\xi}) \, \tilde{w}_k,
\end{align}
where $\tilde{w}_k$ denotes the quadrature weight associated with the quadrature point $\tilde{\xi}_k$. In this work, we employ Gauss quadrature, an integration rule with optimal algebraic accuracy that integrates all polynomials of degree at most $2N^\mathrm{Q}-1$ exactly.

The differentiation matrix $\mathcal{D} \in \mathbb{R}^{({N^\mathrm{d}} + 1) \times ({N^\mathrm{d}} + 1)}$ can be computed directly as
\begin{align}
\label{eq:operator_differentiation_matrix}
\mathcal{D}_{i, j} &\coloneq \frac{\partial \ell_j}{\partial \xi}(\xi_i).
\end{align}

The stiffness matrix with weight $w(\xi)$ is defined as
\begin{align}
\label{eq:operator_stiffness_matrix}
\tilde{\mathcal{S}}_{ij} &\coloneq \int_{-1}^{1} \ell_i \frac{\mathrm{d} \ell_j}{d \xi} w(\xi) \,\mathrm{d}\xi = \left( \tilde{\mathcal{M}} \mathcal{D} \right)_{i,j}.
\end{align}

The so-called ``lifting'' matrix is used to compute surface integrals and is given by
\begin{align}
\label{eq:operator_lifting_matrix}
\mathcal{B}_{ij} :&= \oint_{\partial E} \ell_i(\xi) \ell_j(\xi) w(\xi) \text{\boldmath$n$} \,\mathrm{d}\xi
=
\begin{cases}
- w(-1) & \text{for $i=j=0$,} \\
w(1) & \text{for $i=j={N^\mathrm{d}}$,} \\
0 & \text{else.}
\end{cases}
\end{align}
Here, $\text{\boldmath$n$} \in \{-1, 1\}$ is the outwards pointing unit normal vector in $1$D.

The \gls{dgsem} employs numerical fluxes $f^*\colon \mathbb{R}^2\mapsto \mathbb{R}$ based on the inner and outer values at element interfaces $\xi \in \{-1,1\}$.
This yields an element operator of the form
\begin{align}
\label{eq:operator_numerical_flux_vector}
\underline{f}^* \coloneq \left[f^*|_{\xi=-1}, 0, \dots, 0, f^*|_{\xi=1}\right] \in\mathbb{R}^{N^\mathrm{d}+1}.
\end{align}

% The following summation-by-parts (SBP) property is the discrete equivalent of an integration by parts and will be used to derive the so-called strong form of the DGSEM discretization, which is how the methods are implemented.
% For the two sets of previously defined discrete operators, $\{ \mathcal{D}, \mathcal{B}, \mathcal{M}^{(0, 0)}, \mathcal{S}^{(0, 0)}\}$ and $\{ \mathcal{D}, \mathcal{B}, \mathcal{M}_{LGL}, \mathcal{S}_{LGL}\}$, it holds that
% \begin{align}
%  \mathcal{D} &=  \mathcal{M}^{-1} \  \mathcal{S} \text{ and }  \mathcal{S} +  \mathcal{S}^T =  \mathcal{B}, \label{sbp_property}
% \end{align}
% with $ \mathcal{M}$ and $ \mathcal{S}$ being the mass and stiffness matrices of the respective approach.
% The SBP property~\eqref{sbp_property} was initially shown by \cite{kopriva_quadrature_2010} for LGL and Gau\ss{} integration and also holds for the exact integration operators when $\alpha=\beta=0$.

\subsubsection{DG discretization}
\label{sec:dg_discretization}

Starting point of the discretization is the convection--dispersion--reaction equation for general symmetric column geometries with smoothly varying cross-sectional area defined on $(0,T^\mathrm{end})\times\Omega$
\begin{align}
    \label{eq:generalized_bulk_with_filmDiffusion}
    \frac{\partial c^\mathrm{b}}{\partial t}
    &=
    -\frac{Q}{A(x)}
    \frac{\partial c^\mathrm{b}}{\partial x}
    +
    \frac{1}{A(x)}
    \frac{\partial}{\partial x}
    \left(
    A(x)D^x(x)
    \frac{\partial c^\mathrm{b}}{\partial x}
    \right)
    \nonumber\\
    &\phantom{={}}
    +
    \frac{1-\varepsilon^\mathrm{b}}{\varepsilon^\mathrm{b}}
    \frac{3k^\mathrm{f}(x)}{R^\mathrm{p}}
    \left(
    \left.c^\mathrm{p}\right|_{r=R^\mathrm{p}}
    -
    c^\mathrm{b}
    \right).
\end{align}
First, we partition the spatial domain into $N^\mathrm{e}\in\mathbb{N}$ elements $\Omega_i\coloneq(x_i,x_{i+1})$, such that
\begin{align}
    \overline{\Omega} = \bigcup_i \overline{\Omega_i}.
\end{align}

To improve the numerical conditioning and to enable operator reuse, we map the closure of each element onto the computational reference element $[-1,1]$ via the affine mapping
\begin{align}
    \label{eq:general_reference_mapping}
    \xi(x)
    \coloneq
    2\frac{x-x_i}{\Delta x_i}-1,
    \qquad
    x\in[x_i,x_{i+1}],
\end{align}
with Jacobian
\begin{align}
    \frac{\mathrm d\xi}{\mathrm dx}
    =
    \frac{2}{\Delta x_i},
\end{align}
and inverse
\begin{align}
    x_i^\mathrm{e}(\xi)
    =
    (\xi+1)\frac{\Delta x_i}{2}+x_i,
    \qquad
    \xi\in[-1,1],
\end{align}
where $\Delta x_i\coloneq x_{i+1}-x_i$.

Multiplying \cref{eq:generalized_bulk_with_filmDiffusion} by a smooth test function
$\phi\in\mathcal C^\infty(\Omega)$
and integrating over an element gives
\begin{align*}
    \int_{\Omega_i}
    \phi
    \frac{\partial c^\mathrm{b}}{\partial t}
    A(x)
    \,\mathrm dx
    &=
    \int_{\Omega_i}
    \phi
    \Bigg(
    -Q
    \frac{\partial c^\mathrm{b}}{\partial x}
    +
    \frac{\partial}{\partial x}
    \left(
    A(x)D^x(x)
    \frac{\partial c^\mathrm{b}}{\partial x}
    \right)
    \\
    &\qquad\qquad
    +
    \frac{1-\varepsilon^\mathrm b}{\varepsilon^\mathrm b}
    \frac{3k^\mathrm f(x)}{R^\mathrm p}
    A(x)
    \left(
    \left.c^\mathrm p\right|_{r=R^\mathrm p}
    -c^\mathrm{b}
    \right)
    \Bigg)
    \,\mathrm dx.
\end{align*}

We transform the equation to the reference element and introduce the shorthands
\begin{align*}
    w^{AD^\mathrm{x}}(\xi) \coloneq A(x_i^\mathrm e(\xi)) D^x(x_i^\mathrm e(\xi)),
    \quad
    A(\xi) \coloneq A(x_i^\mathrm e(\xi)),
\end{align*}
which yields
\begin{align*}
    \int_{-1}^1
    \phi
    \frac{\partial c^\mathrm{b}}{\partial t}
    A(\xi)
    \frac{\Delta x_i}{2}
    \,\mathrm d\xi
    &=
    \int_{-1}^1
    \phi
    \left(
    -
    Q
    \frac{\partial c^\mathrm{b}}{\partial\xi}
    \frac{2}{\Delta x_i}
    +
    \frac{\partial}{\partial\xi}
    \frac{2}{\Delta x_i}
    \left(
    w
    \frac{\partial c^\mathrm{b}}{\partial\xi}
    \frac{2}{\Delta x_i}
    \right)
    \right)
    \,\mathrm d\xi
    \\
    &\phantom{={}}
    +
    \int_{-1}^1
    \phi
    \frac{1-\varepsilon^\mathrm b}{\varepsilon^\mathrm b}
    \frac{3k^\mathrm f(x_i^\mathrm e(\xi))}{R^\mathrm p}
    A(\xi)
    \left(
    \left.c^\mathrm p\right|_{r=R^\mathrm p}
    -
    c^\mathrm{b}
    \right)
    \frac{\Delta x_i}{2}
    \,\mathrm d\xi.
\end{align*}

The Jacobian of the integral substitution cancels and we perform an integration by parts to obtain
\begin{align*}
    \int_{-1}^1 \phi \frac{\partial c^\mathrm{b}}{\partial t} A(\xi) \,\mathrm d\xi
    &=
    -
    \frac{2}{\Delta x_i}
    \left[
    \phi
    \left(
    Qc^\mathrm{b}
    -
    wg
    \right)
    \right]_{-1}^{1}
    \\
    &\phantom{={}}
    +
    \int_{-1}^1
    \frac{\partial\phi}{\partial\xi}
    \left(
    Q c^\mathrm{b}
    \frac{2}{\Delta x_i}
    -
    \frac{2}{\Delta x_i}
    wg
    \right)
    \,\mathrm d\xi
    \\
    &\phantom{={}}
    +
    \int_{-1}^1
    \phi
    \frac{1-\varepsilon^\mathrm b}{\varepsilon^\mathrm b}
    \frac{3k^\mathrm f(x_i^\mathrm e(\xi))}{R^\mathrm p}
    A(\xi)
    \left(
    \left.c^\mathrm p\right|_{r=R^\mathrm p}
    -
    c^\mathrm{b}
    \right)
    \,\mathrm d\xi.
\end{align*}
Here, we followed the approach of \cite{bassi_high-order_1997} and introduced the auxiliary variable
\begin{align}
   g \coloneq \frac{2}{\Delta x_i} \frac{\partial c^\mathrm{b}}{\partial\xi},
\end{align}
to rewrite the equation as a first order system.
The auxiliary weak formulation reads
\begin{align*}
    \int_{-1}^1
    \phi g
    \,\mathrm d\xi
    &=
    \frac{2}{\Delta x_i}
    \int_{-1}^1
    \phi
    \frac{\partial c^\mathrm{b}}{\partial\xi}
    \,\mathrm d\xi
    \\
    &=
    \frac{2}{\Delta x_i}
    \left[
    \phi c^\mathrm{b}
    \right]_{-1}^{1}
    -
    \frac{2}{\Delta x_i}
    \int_{-1}^1
    \frac{\partial\phi}{\partial\xi}
    c^\mathrm{b}
    \,\mathrm d\xi.
\end{align*}

We note that a solution to the weak formulation is not necessarily a solution to the original equations and vice versa. However, a sufficiently smooth solution satisfies both the original and weak form.

Next, we seek approximate solutions of the weak form in the polynomial space. To this end, we employ interpolating Lagrange polynomials of degree $N^\mathrm{d}$ on the computational element and obtain for $t\in(0,T^\mathrm{end}]$ and $\xi\in[-1,1]$
\begin{subequations}
\begin{align}
    c^\mathrm{b}(t,\xi)
    &\approx
    \sum_{k=0}^{N^\mathrm d}
    \underline{c}^\mathrm b_k(t)
    \ell_k(\xi)
    =:c_h,
    \\
    g(t,\xi)
    &\approx
    \sum_{k=0}^{N^\mathrm d}
    \underline{g}_k(t)
    \ell_k(\xi)
    =:g_h,
    \\
    c^\mathrm p(t,\xi)
    &\approx
    \sum_{k=0}^{N^\mathrm d}
    \underline{c}^\mathrm p_k(t)
    \ell_k(\xi)
    =:c_h^\mathrm p.
\end{align}
\end{subequations}

Here, $\underline{c}^\mathrm{b},\underline{g},\underline{c}^\mathrm{p}$ denote the vectors of polynomial coefficients.
Discretizing the continuous weak form, we reduce the infinite-dimensional test space $\mathcal C^\infty([-1,1])$ to the finite-dimensional polynomial space $\mathbb P^{N^\mathrm d}([-1,1])$.
Choosing the Lagrange basis functions $\ell_j\in\mathcal L^{N^\mathrm d}$ as test functions yields the discrete weak formulation
\begin{align*}
    \int_{-1}^1
    \ell_j
    \frac{\partial c_h}{\partial t}
    A(\xi)
    \,\mathrm d\xi
    &=
    -\frac{2}{\Delta x_i}
    \left[
    \ell_j
    \left(
    Qc_h^{*,\mathrm a}
    -
    w^{AD^\mathrm{x}} g_h^{*,\mathrm d}
    \right)
    \right]_{-1}^{1}
    \\
    &\phantom{={}}
    +
    \int_{-1}^1
    \frac{2}{\Delta x_i}
    \frac{\partial\ell_j}{\partial\xi}
    \left(
    Qc_h
    -
    w^{AD^\mathrm{x}} g_h
    \right)
    \,\mathrm d\xi
    \\
    &\phantom{={}}
    +
    \int_{-1}^1
    \ell_j
    \frac{1-\varepsilon^\mathrm b}{\varepsilon^\mathrm b}
    \frac{3k^\mathrm f(x_i^\mathrm e(\xi))}{R^\mathrm p}
    A(\xi)
    \left(
    \left.c^\mathrm p\right|_{r=R^\mathrm p}
    -
    c_h
    \right)
    \,\mathrm d\xi,
    \\
    \int_{-1}^1
    \ell_j g_h
    \,\mathrm d\xi
    &=
    \frac{2}{\Delta x_i}
    \left[
    \ell_j
    c_h^{*,\mathrm d}
    \right]_{-1}^{1}
    -
    \frac{2}{\Delta x_i}
    \int_{-1}^1
    \frac{\partial\ell_j}{\partial\xi}
    c_h
    \,\mathrm d\xi.
\end{align*}

Collecting the element integrals into the discrete operators introduced in \cref{sec:nodal_element_toolbox}, we obtain the element-wise matrix formulation
\begin{align}
    \label{eq:generalized_matrix_operator_dg_discretization}
    \mathcal{M}^{\mathrm{A}}
    \frac{\partial\underline{c}^\mathrm{b}}{\partial t}
    &=
    -\frac{2}{\Delta x_i}
    \left(
    \mathcal{B} \, Q \, \underline c^{*,\mathrm a}
    -
    \mathcal{B}^{AD^\mathrm{x}} \underline{g}^{*,\mathrm d}
    \right)
    +
    \frac{2}{\Delta x_i}
    \left(
    \mathcal{D}^\mathrm{T}\mathcal{M}^{(0,0)}
    \, Q \, \underline{c}^\mathrm{b}
    -
    \tilde{\mathcal{S}}^{AD^\mathrm{x}} \underline{g}
    \right)
    \nonumber \\
    &\phantom{={}}
    +
    \frac{1-\varepsilon^\mathrm b}{\varepsilon^\mathrm b}
    \frac{3}{R^\mathrm p}
    \tilde{\mathcal{M}}^{Ak^\mathrm{f}}
    \left(
    \underline c^\mathrm p
    -
    \underline{c}^\mathrm{b}
    \right),
    \\
    \label{eq:auxiliary_matrix_operator_dg_discretization}
    \mathcal{M}^{(0,0)}
    \underline{g}
    &=
    \frac{2}{\Delta x_i}
    \left(
    \mathcal{B}\underline c^{*,\mathrm d}
    -
    \mathcal{D}^\mathrm{T}\mathcal{M}^{(0,0)}
    \underline{c}^\mathrm{b}
    \right).
\end{align}
The vectors $\underline c^{*,\mathrm a}$, $\underline c^{*,\mathrm d}$, and $\underline{g}^{*,\mathrm d}$ are the numerical flux operators as defined in \cref{eq:operator_numerical_flux_vector}.
We recover the element-wise matrix operators used in \cref{eq:generalized_matrix_operator_dg_discretization} for both the radial flow and frustum geometry (see Model~\eqref{model:radial_flow_bulk} and~\eqref{model:frustum_bulk}) by substituting the corresponding cross-sectional area function $A(x)$ into the weighted operators, in accordance with the preparations made in \cref{sec:nodal_element_toolbox}.
The cross-sectional area is given by
\begin{align}
    A(x)=
    \begin{cases}
    2\pi L\,x,
    &
    \text{radial flow},
    \\[1ex]
    \displaystyle
    \pi
    \left(
    r_0+\frac{x}{H}(r_H-r_0)
    \right)^2,
    &
    \text{frustum flow}.
    \end{cases}
\end{align}

The weighted stiffness matrix $\tilde{\mathcal{S}}^{AD^\mathrm{x}}$ involves a spatially varying dispersion parameter, which we do not know in general, and thus requires numerical integration and is defined as per \eqref{eq:operator_stiffness_matrix} with weight $A(x^\mathrm{e}_i(\xi)) D^\mathrm{x}(x^\mathrm{e}_i(\xi))$.
Similarly, we define the weighted mass matrix $\tilde{\mathcal{M}}^{Ak^\mathrm{f}}$ per \eqref{eq:operator_gauss_mass_matrix} with weight $A(x^\mathrm{e}_i(\xi)) k^\mathrm{f}(x^\mathrm{e}_i(\xi))$.

The mass matrix $\mathcal{M}^{\mathrm{A}}$ can be computed exactly as a linear combination of Jacobi-weighted mass matrices~\eqref{eq:operator_exact_jacobi_mass_matrix}.
For the radial flow geometry, we get
\begin{subequations}
\begin{align}
    \mathcal{M}^{\mathrm{A}}
    &=
    2\pi L
    \left(
    x_i\,\mathcal{M}^{(0,0)}
    +
    \frac{\Delta x_i}{2}\,
    \mathcal{M}^{(0,1)}
    \right).
\intertext{For the frustum flow geometry, expanding the quadratic area function gives}
    \mathcal{M}^{\mathrm A}
    &=
    \pi
    \Biggl[
    \left(
    r_0^2
    +\frac{2r_0x_i}{H}(r_H-r_0)
    +\frac{x_i^2}{H^2}(r_H-r_0)^2
    \right)
    \mathcal{M}^{(0,0)}
    \nonumber\\
    &\phantom{=\pi}
    +
    \left(
    \frac{2r_0(r_H-r_0)}{H}\frac{\Delta x_i}{2}
    +\frac{(r_H-r_0)^2}{H^2} x_i \Delta x_i
    \right)
    \mathcal{M}^{(0,1)}
    \nonumber\\
    &\phantom{=\pi}
    +
    \frac{(r_H-r_0)^2}{H^2}
    \left(\frac{\Delta x_i}{2}\right)^2
    \mathcal{M}^{(0,2)}
    \Biggr].
\end{align}
\end{subequations}

The lifting matrix $\mathcal{B}_{AD^\mathrm{x}}$ is defined via \eqref{eq:operator_lifting_matrix} with weight $A(x^\mathrm{e}_i(\xi)) D^\mathrm{x}(x^\mathrm{e}_i(\xi))$.

We further simplify the auxiliary \cref{eq:auxiliary_matrix_operator_dg_discretization} using the summation-by-parts property \citep[see][]{Kopriva2010}, which holds in the non-weighted space, yielding
\begin{align}
    \label{eq:auxiliary_strong_form}
    \underline{g}
    =
    \frac{2}{\Delta x_i}
    \left( \left(\mathcal{M}^{(0,0)}\right)^{-1} B
        \left( \underline c^{*,\mathrm d} - \underline{c}^\mathrm{b} \right)
        + D \underline{c}^\mathrm{b}
    \right).
\end{align}

The complementing particle model equations defined in \cref{sec:particle_models} are discretized following \cite{BREUER2023108340} which, for completeness, is summarized in the Appendix~\ref{appendix:sec:particle_discretization}.


\subsubsection{Numerical Fluxes}

For the linear advection term, we approximate the interface fluxes with an upwind numerical flux for $v>0$
\begin{subequations}
\label{eq:radial_flow_numerical_fluxes}
\begin{align}
c^{*,\mathrm{a}}_h \left( u (c_h)^-, u (c_h)^+ \right)
&\coloneq u (c_h)^-,
\intertext{where the superscripts $-$ and $+$ denote the inner and outer interface value respectively.
For the dispersion interface fluxes, we employ a central numerical flux \citep{bassi_high-order_1997}}
g^{*,\mathrm{d}}_h \left( g_h^-, g_h^+ \right)
&\coloneq \frac{g_h^+ + g_h^-}{2},
\\
c^{*,\mathrm{d}}_h \left( c_h^-, c_h^+ \right)
&\coloneq \frac{c_h^+ + c_h^-}{2}.
\end{align}
\end{subequations}


\subsubsection{Boundary Conditions}

Danckwerts boundary conditions \eqref{eq:danckwerts_bc} are imposed weakly through the numerical fluxes
\begin{subequations}
\begin{align}
c^{a,*}+g^{d,*}\big|_{\rho=P^\mathrm{c}} &\coloneq u c^b_\text{in}, &\quad c^{a,*}+g^{d,*}\big|_{\rho=P} &\coloneq u (c_h)^-.
\intertext{
The model does not prescribe boundary conditions for the auxiliary equations.
We hence define the numerical flux for inflow and outflow boundaries as \citep{bassi_high-order_1997}
}
c^{d,*}\big|_{\rho=P^\mathrm{c}} &\coloneq c^b\big|_{\rho=P^\mathrm{c}}, &\quad c^{d,*}\big|_{\rho=P} &\coloneq c^b\big|_{\rho=P}.
\end{align}
\end{subequations}